\section*{Capítulo \ref{ch:g10:common-ancestry} --- Planos de organização e ancestralidade comum}

\begin{solution}{exo:g10:common-ancestry:1}
Simetria bilateral com uma extremidade anterior e uma posterior; um
esqueleto interno com coluna vertebral; um cordão nervoso dorsal acima
dela; um tubo digestório ventral com o coração por baixo dele, perto da
frente (e membros pares ou nadadeiras).
\end{solution}

\begin{solution}{exo:g10:common-ancestry:2}
Órgãos com a mesma estrutura e a mesma posição no \omterm{def:g10:common-ancestry:bodyplan}{plano de organização},
seja qual for a função deles. O membro anterior dos tetrápodes: úmero,
rádio e ulna, ossos do carpo, dedos --- num braço, numa asa, numa
nadadeira.
\end{solution}

\begin{solution}{exo:g10:common-ancestry:3}
Análogas: mesma função (o voo), estrutura inteiramente diferente ---
ossos, músculos e pele na ave, uma membrana do esqueleto externo na
borboleta. Elas não dizem nada sobre parentesco.
\end{solution}

\begin{solution}{exo:g10:common-ancestry:4}
O lagarto: a menor caixa que contém rã e lagarto é a dos
``tetrápodes''; a truta só está na caixa maior, a dos ``\omterm{def:g10:common-ancestry:bodyplan}{vertebrados}''.
\end{solution}

\begin{solution}{exo:g10:common-ancestry:5}
Fósseis com combinações intermediárias de caracteres; os estágios
compartilhados dos embriões de \omterm{def:g10:common-ancestry:bodyplan}{vertebrados}; a concordância das
comparações moleculares de sequência com os agrupamentos anatômicos.
\end{solution}

\begin{solution}{exo:g10:common-ancestry:6}
Aos mamíferos (vestígios de pelos, leite; membro de tetrápode;
respiração aérea e, nos mamíferos, um desenvolvimento amniótico). Os
ossos do quadril são o resto dos membros posteriores: as baleias
descendem de mamíferos terrestres de quatro patas.
\end{solution}

\begin{solution}{exo:g10:common-ancestry:7}
O crocodilo entra em ``amniotas'' com o lagarto e o pombo (e não em
``mamíferos''); a salamandra entra em ``tetrápodes'' fora de
``amniotas'', junto com a rã.
\end{solution}

\begin{solution}{exo:g10:common-ancestry:8}
Um grupo de parentesco tem de conter todos os descendentes do ancestral
comum dele. O ancestral comum à truta e ao peixe pulmonado também é o
ancestral da rã e de todo tetrápode; um grupo que contenha a truta e o
peixe pulmonado mas exclua a rã deixa de fora alguns descendentes.
``Peixes'' nomeia um modo de vida, não uma linhagem.
\end{solution}

\begin{solution}{exo:g10:common-ancestry:9}
Chimpanzé (0), camundongo (9), galinha (13), rã (18), atum (21),
levedura (45): mamífero, amniota, tetrápode, \omterm{def:g10:common-ancestry:bodyplan}{vertebrado}, eucarionte ---
as caixas em ordem.
\end{solution}

\begin{solution}{exo:g10:common-ancestry:10}
Estruturas mantidas no peixe adulto aparecem transitoriamente nos
nossos embriões e são remodeladas: as bolsas branquiais viram partes do
ouvido, da garganta e de glândulas, e a cauda é reabsorvida. Um embrião
que percorre um estágio ancestral é o que a herança a partir de um
ancestral comum prevê e o que um projeto para uma finalidade não prevê.
\end{solution}

\begin{solution}{exo:g10:common-ancestry:11}
Se as aves descendem de ancestrais semelhantes a répteis, algum
ancestral tem de ter tido caracteres de réptil e os primeiros
caracteres de ave ao mesmo tempo; encontrar um datado entre os dois
grupos é uma previsão confirmada. Sem a ancestralidade comum, uma
criatura que mistura os traços de dois tipos separados não tem razão
para existir.
\end{solution}

\begin{solution}{exo:g10:common-ancestry:12}
Análogos: mesma função e projeto geral parecido, mas origem embrionária
diferente e estrutura invertida. A \omterm{def:g10:common-ancestry:homology}{homologia} implicaria que o ancestral
comum dos moluscos e dos \omterm{def:g10:common-ancestry:bodyplan}{vertebrados} já tinha um olho de câmara; a
analogia implica que ele tinha no máximo uma simples mancha sensível à
luz, e que as duas linhagens construíram as câmaras separadamente.
\end{solution}

\begin{solution}{exo:g10:common-ancestry:13}
O ser humano e o camundongo compartilham um ancestral comum mais
recente do que o que cada um compartilha com a galinha; as duas
linhagens estão separadas da linhagem da galinha há o mesmo tempo, de
modo que as duas acumularam mais ou menos o mesmo número de mudanças em
relação a ela. As distâncias humano--galinha e camundongo--galinha
medem o mesmo intervalo.
\end{solution}

\begin{solution}{exo:g10:common-ancestry:14}
As serpentes têm o ovo amniótico, o esqueleto e o crânio dos lagartos,
e algumas (jiboias, pítons) conservam vestígios dos membros posteriores
e da pelve; existem serpentes fósseis com perninhas. ``Tetrápode''
quer dizer descendente do ancestral de quatro membros, tendo ou não
conservado os membros.
\end{solution}

\begin{solution}{exo:g10:common-ancestry:15}
O modo de vida explica a forma de uma nadadeira, mas não por que a
nadadeira de uma baleia contém os ossos de um braço arrumados como os
de um morcego --- uma pá achatada não precisa de punho. Ele não explica
as bolsas branquiais num embrião humano, que não tem uso para elas. E
não explica por que uma \omterm{def:g10:chemistry-of-life:families}{proteína} respiratória, invisível ao ambiente,
difere entre as \omterm{def:g10:biodiversity-scales:species}{espécies} exatamente na medida que os esqueletos
preveem. Só a herança a partir de ancestrais comuns dá conta das três
coisas.
\end{solution}

\begin{solution}{pb:g10:common-ancestry:1}
\textbf{1.} A coluna vertebral: os \omterm{def:g10:common-ancestry:bodyplan}{vertebrados}.

\textbf{2.} Quatro membros define rã, lagarto, pombo e camundongo: os
tetrápodes. O ovo amniótico define lagarto, pombo e camundongo: os
amniotas.

\textbf{3.} Não: um caráter presente numa só \omterm{def:g10:biodiversity-scales:species}{espécie} não agrupa nada.
Ele identifica a \omterm{def:g10:biodiversity-scales:species}{espécie} (e definiria um grupo se mais parentes ---
outras aves, outros mamíferos --- fossem acrescentados).

\textbf{4.} \omterm{def:g10:common-ancestry:bodyplan}{Vertebrados} $\supset$ tetrápodes $\supset$ amniotas, com a
truta fora da segunda caixa e a rã fora da terceira; o lagarto, o pombo
e o camundongo dentro da terceira.

\textbf{5.} A da rã: mesmos ossos, mesma ordem. Corresponde ao caráter
``quatro membros'', que a nadadeira não tem.

\textbf{6.} Lagarto e pombo (12). A truta (19--21 de todas).

\textbf{7.} Pombo (13), lagarto (14), rã (18), truta (21).

\textbf{8.} Lagarto (16), pombo (17), camundongo (18), truta (19): os
amniotas a distâncias mais ou menos iguais, a truta mais além ---
coerente com as caixas (a rã está mais perto de todo amniota do que da
truta).

\textbf{9.} Que o lagarto e o pombo se separaram um do outro depois de
a linhagem deles ter se separado da do camundongo: dentro dos
amniotas, lagarto e pombo são parentes mais próximos. A tabela não
consegue dizer: nenhum caráter dela é compartilhado só pelo lagarto e
pelo pombo.

\textbf{10.} A linhagem da truta se separou da linhagem comum das
outras quatro antes que estas se separassem entre si: cada uma está
apartada da truta há o mesmo tempo, de modo que cada uma acumulou mais
ou menos o mesmo número de diferenças.

\textbf{11.} A truta fora de todas as outras; os quatro tetrápodes
juntos; os três amniotas juntos dentro deles.

\textbf{12.} A anatomia e uma \omterm{def:g10:chemistry-of-life:families}{proteína} respiratória não têm razão para
concordar a menos que ambas registrem a mesma história. Duas
testemunhas independentes que contam a mesma história tornam uma
semelhança acidental implausível.

\textbf{13.} Em ``mamíferos'', ao lado do camundongo (pelos, leite e a
menor distância molecular). A asa dele é um membro anterior de
tetrápode modificado, análogo ao do pombo, e não homólogo ao voo
emplumado; a molécula (13 do pombo, 8 do camundongo) confirma.

\textbf{14.} O citocromo c muda lentamente demais para separar parentes
próximos, e não é o \omterm{def:g10:universal-dna:gene}{gene} responsável pelas diferenças visíveis. Ele
informa sobre parentesco profundo, e não sobre os \omterm{def:g10:universal-dna:gene}{genes} que fazem de
uma \omterm{def:g10:biodiversity-scales:species}{espécie} o que ela é.

\textbf{15.} Os eucariontes: \omterm{def:g10:cells-common-unit:cell}{células} com \omterm{def:g10:cells-common-unit:organelle}{núcleo} e \omterm{def:g10:cells-common-unit:organelle}{mitocôndrias}, que
compartilham essa \omterm{def:g10:chemistry-of-life:families}{proteína} respiratória.

\textbf{16.} As 13 diferenças foram acumuladas por duas linhagens ao
longo de 320 milhões de anos: $13/2 \approx 6.5$ por linhagem, ou seja
cerca de 2 por 100 milhões de anos ao longo de uma linhagem.

\textbf{17.} Rã--camundongo 18, rã--pombo 17: média 17.5, ou seja 8.75
por linhagem; a 2 por 100 milhões de anos, cerca de 430 milhões de anos.

\textbf{18.} A truta dá média de cerca de 20: 10 por linhagem, cerca de
500 milhões de anos.

\textbf{19.} As estimativas (430, 500) são mais antigas que as datas
dos fósseis (370, mais de 420), mas da ordem certa e na sequência
certa; a taxa não é exatamente constante e as contagens estão
arredondadas, de modo que uma concordância dentro de 20\% é o melhor
que o método permite.

\textbf{20.} Os dois estabelecem o mesmo parentesco encaixado ---
truta, depois rã, depois os três amniotas ---, enquanto a molécula
acrescenta um relógio: uma estimativa de há quanto tempo cada ramo se
separou.
\end{solution}
